\cleardoublepage
\phantomsection
\addcontentsline{toc}{chapter}{Nombres entiers \& divisibilité}

\begin{exercice}{1}{Divisibilité}

Complète les tableaux suivants:
\bigskip

\textbf{A.}
\medskip

\def\arraystretch{1.2}
\begin{tabularx}{\linewidth}{|c|*{11}{>{\centering\arraybackslash}X|}}
\hline
\textbf{divisible par} & 2 & 3 & 4 & 5 & 6 &  9 &  10 &  25 & 50 &  75 &  100 \\\hline
13464& & & & & & & & & & & \\\hline
1800& & & & & & & & & & & \\\hline
36275& & & & & & & & & & & \\\hline
41080& & & & & & & & & & & \\\hline
410205& & & & & & & & & & & \\\hline
12350& & & & & & & & & & & \\\hline
\end{tabularx}
\vfill

\textbf{B.}
\medskip

\begin{tabularx}{\linewidth}{|c|*{11}{>{\centering\arraybackslash}X|}}
\hline
\textbf{divisible par} & 2 & 3 & 4 & 5 & 6 &  9 &  10 &  25 & 50 &  75 &  100 \\\hline
43464& & & & & & & & & & & \\\hline
18800& & & & & & & & & & & \\\hline
6275& & & & & & & & & & & \\\hline
91080& & & & & & & & & & & \\\hline
610290& & & & & & & & & & & \\\hline
52350& & & & & & & & & & & \\\hline
\end{tabularx}
\vfill

\textbf{C.}
\medskip

\begin{tabularx}{\linewidth}{|c|*{11}{>{\centering\arraybackslash}X|}}
\hline
\textbf{divisible par} & 2 & 3 & 4 & 5 & 6 &  9 &  10 &  25 & 50 &  75 &  100 \\\hline
17820& & & & & & & & & & & \\\hline
46200& & & & & & & & & & & \\\hline
71180 & & & & & & & & & & & \\\hline
610215& & & & & & & & & & & \\\hline
62352& & & & & & & & & & & \\\hline
256400& & & & & & & & & & & \\\hline
\end{tabularx}
\vfill

\textbf{D.}
\medskip

\begin{tabularx}{\linewidth}{|c|*{11}{>{\centering\arraybackslash}X|}}
\hline
\textbf{divisible par} & 2 & 3 & 4 & 5 & 6 &  9 &  10 &  25 & 50 &  75 &  100 \\\hline
1880& & & & & & & & & & & \\\hline
28530& & & & & & & & & & & \\\hline
925& & & & & & & & & & & \\\hline
275985& & & & & & & & & & & \\\hline
15\_\_  \_\_& x & & x & x & & & x & & & & \\\hline
387\_\_  \_\_& x & x & x & & x & x & & & & & \\\hline
\end{tabularx}

\end{exercice}
\begin{solution}{1}

\textbf{A.}
\bigskip

\begin{tabular}{|l|l|l|l|l|l|l|l|l|l|l|l|}
\hline
\textbf{divisible par...} & \textbf{2} & \textbf{3} & \textbf{4} & \textbf{5} & \textbf{6} & \textbf{9} & \textbf{10} & \textbf{25} & \textbf{50} & \textbf{75} & \textbf{100} \\
\hline
\textbf{13464}  & $\checkmark$ & $\checkmark$ & $\checkmark$ & $\times$ & $\checkmark$ & $\checkmark$ & $\times$ & $\times$ & $\times$ & $\times$ & $\times$ \\
\hline
\textbf{1800}   & $\checkmark$ & $\checkmark$ & $\checkmark$ & $\checkmark$ & $\checkmark$ & $\checkmark$ & $\checkmark$ & $\checkmark$ & $\checkmark$ & $\checkmark$ & $\checkmark$ \\
\hline
\textbf{36275}  & $\times$ & $\times$ & $\times$ & $\checkmark$ & $\times$ & $\times$ & $\times$ & $\checkmark$ & $\times$ & $\times$ & $\times$ \\
\hline
\textbf{41080}  & $\checkmark$ & $\times$ & $\checkmark$ & $\checkmark$ & $\times$ & $\times$ & $\checkmark$ & $\times$ & $\times$ & $\times$ & $\times$ \\
\hline
\textbf{410205} & $\times$ & $\checkmark$ & $\times$ & $\checkmark$ & $\times$ & $\times$ & $\times$ & $\times$ & $\times$ & $\times$ & $\times$ \\
\hline
\textbf{12350}  & $\checkmark$ & $\times$ & $\times$ & $\checkmark$ & $\times$ & $\times$ & $\checkmark$ & $\checkmark$ & $\checkmark$ & $\times$ & $\times$ \\
\hline
\end{tabular}
\bigskip

\textit{Rappel des critères : somme des chiffres pour 3 et 9 ; deux derniers chiffres pour 4, 25, 50, 100 ; dernier chiffre pour 2, 5, 10 ; 6 = 2 et 3 ; 75 = 25 et 3.}
\bigskip

\textbf{B.}
\bigskip

\begin{tabular}{|l|l|l|l|l|l|l|l|l|l|l|l|}
\hline
\textbf{divisible par...} & \textbf{2} & \textbf{3} & \textbf{4} & \textbf{5} & \textbf{6} & \textbf{9} & \textbf{10} & \textbf{25} & \textbf{50} & \textbf{75} & \textbf{100} \\
\hline
\textbf{43464}  & $\checkmark$ & $\checkmark$ & $\checkmark$ & $\times$ & $\checkmark$ & $\times$ & $\times$ & $\times$ & $\times$ & $\times$ & $\times$ \\
\hline
\textbf{18800}  & $\checkmark$ & $\times$ & $\checkmark$ & $\checkmark$ & $\times$ & $\times$ & $\checkmark$ & $\checkmark$ & $\checkmark$ & $\times$ & $\checkmark$ \\
\hline
\textbf{6275}   & $\times$ & $\times$ & $\times$ & $\checkmark$ & $\times$ & $\times$ & $\times$ & $\checkmark$ & $\times$ & $\times$ & $\times$ \\
\hline
\textbf{91080}  & $\checkmark$ & $\checkmark$ & $\checkmark$ & $\checkmark$ & $\checkmark$ & $\checkmark$ & $\checkmark$ & $\times$ & $\times$ & $\times$ & $\times$ \\
\hline
\textbf{610290} & $\checkmark$ & $\checkmark$ & $\times$ & $\checkmark$ & $\checkmark$ & $\checkmark$ & $\checkmark$ & $\times$ & $\times$ & $\times$ & $\times$ \\
\hline
\textbf{52350}  & $\checkmark$ & $\checkmark$ & $\times$ & $\checkmark$ & $\checkmark$ & $\times$ & $\checkmark$ & $\checkmark$ & $\checkmark$ & $\checkmark$ & $\times$ \\
\hline
\end{tabular}
\bigskip

\textbf{C.}
\bigskip

\begin{tabular}{|l|l|l|l|l|l|l|l|l|l|l|l|}
\hline
\textbf{divisible par...} & \textbf{2} & \textbf{3} & \textbf{4} & \textbf{5} & \textbf{6} & \textbf{9} & \textbf{10} & \textbf{25} & \textbf{50} & \textbf{75} & \textbf{100} \\
\hline
\textbf{17820}  & $\checkmark$ & $\checkmark$ & $\checkmark$ & $\checkmark$ & $\checkmark$ & $\checkmark$ & $\checkmark$ & $\times$ & $\times$ & $\times$ & $\times$ \\
\hline
\textbf{46200}  & $\checkmark$ & $\checkmark$ & $\checkmark$ & $\checkmark$ & $\checkmark$ & $\times$ & $\checkmark$ & $\checkmark$ & $\checkmark$ & $\checkmark$ & $\checkmark$ \\
\hline
\textbf{71180}  & $\checkmark$ & $\times$ & $\checkmark$ & $\checkmark$ & $\times$ & $\times$ & $\checkmark$ & $\times$ & $\times$ & $\times$ & $\times$ \\
\hline
\textbf{610215} & $\times$ & $\checkmark$ & $\times$ & $\checkmark$ & $\times$ & $\times$ & $\times$ & $\times$ & $\times$ & $\times$ & $\times$ \\
\hline
\textbf{62352}  & $\checkmark$ & $\checkmark$ & $\checkmark$ & $\times$ & $\checkmark$ & $\checkmark$ & $\times$ & $\times$ & $\times$ & $\times$ & $\times$ \\
\hline
\textbf{256400} & $\checkmark$ & $\times$ & $\checkmark$ & $\checkmark$ & $\times$ & $\times$ & $\checkmark$ & $\checkmark$ & $\checkmark$ & $\times$ & $\checkmark$ \\
\hline
\end{tabular}
\bigskip

\textbf{D.}
\bigskip

\begin{tabular}{|l|l|l|l|l|l|l|l|l|l|l|l|}
\hline
\textbf{divisible par...} & \textbf{2} & \textbf{3} & \textbf{4} & \textbf{5} & \textbf{6} & \textbf{9} & \textbf{10} & \textbf{25} & \textbf{50} & \textbf{75} & \textbf{100} \\
\hline
\textbf{1880}   & $\checkmark$ & $\times$ & $\checkmark$ & $\checkmark$ & $\times$ & $\times$ & $\checkmark$ & $\times$ & $\times$ & $\times$ & $\times$ \\
\hline
\textbf{28530}  & $\checkmark$ & $\checkmark$ & $\times$ & $\checkmark$ & $\checkmark$ & $\checkmark$ & $\checkmark$ & $\times$ & $\times$ & $\times$ & $\times$ \\
\hline
\textbf{925}    & $\times$ & $\times$ & $\times$ & $\checkmark$ & $\times$ & $\times$ & $\times$ & $\checkmark$ & $\times$ & $\times$ & $\times$ \\
\hline
\textbf{15280}  & $\checkmark$ & $\times$ & $\checkmark$ & $\checkmark$ & $\times$ & $\times$ & $\checkmark$ & $\times$ & $\times$ & $\times$ & $\times$ \\
\hline
\textbf{275985} & $\times$ & $\checkmark$ & $\times$ & $\checkmark$ & $\times$ & $\checkmark$ & $\times$ & $\times$ & $\times$ & $\times$ & $\times$ \\
\hline
\textbf{38736}  & $\checkmark$ & $\checkmark$ & $\checkmark$ & $\times$ & $\checkmark$ & $\checkmark$ & $\times$ & $\times$ & $\times$ & $\times$ & $\times$ \\
\hline
\end{tabular}
\def\arraystretch{1}

\bigskip
\textit{Pour les lignes à compléter de D : toute réponse correctement justifiée est acceptée. Exemples donnés : \textbf{15280} (divisible par 2, 4, 5, 10 ; somme des chiffres = 16, non divisible par 3) et \textbf{38736} (somme des chiffres = 27, divisible par 3 et 9 ; deux derniers chiffres 36, divisible par 4).}

\end{solution}

\begin{exercice}{2}{Décomposition en facteurs premiers}

\begin{mdframed}[backgroundcolor=vlmgray, innerleftmargin=10pt, innerrightmargin=10pt, innertopmargin=6pt, innerbottommargin=6pt, skipabove=-5pt, skipbelow=5pt]
\textbf{Théorème fondamental de l'arithmétique :}

Chaque entier supérieur à 1 est premier ou un produit de nombres premiers ; la factorisation est unique, mis à part l'ordre des facteurs.
\end{mdframed}


\medskip
Décompose chaque nombre ci-dessous en un produit de facteurs premiers.

\medskip
\def\arraystretch{1.7}
\begin{tabularx}{\linewidth}{*{3}{X}}
\textbf{A.} & \textbf{B.} & \textbf{C.} \\
$255  = \dotfill$ & $1340 = \dotfill$ & $720  = \dotfill$ \\
$450  = \dotfill$ & $140  = \dotfill$ & $378  = \dotfill$ \\
$84   = \dotfill$ & $342  = \dotfill$ & $270  = \dotfill$ \\
$42   = \dotfill$ & $280  = \dotfill$ & $210  = \dotfill$ \\
$600  = \dotfill$ & $144  = \dotfill$ & $117  = \dotfill$ \\
$52   = \dotfill$ & $1450 = \dotfill$ & $480  = \dotfill$ \\
\end{tabularx}

\end{exercice}
\begin{solution}{2}

\begin{tabular}{p{5cm} p{5cm} p{5cm}}
\textbf{A.} & \textbf{B.} & \textbf{C.} \\[6pt]
$255  = 3 \times 5 \times 17$          & $1340 = 2^2 \times 5 \times 67$  & $720  = 2^4 \times 3^2 \times 5$   \\[8pt]
$450  = 2 \times 3^2 \times 5^2$       & $140  = 2^2 \times 5 \times 7$   & $378  = 2 \times 3^3 \times 7$     \\[8pt]
$84   = 2^2 \times 3 \times 7$         & $342  = 2 \times 3^2 \times 19$  & $270  = 2 \times 3^3 \times 5$     \\[8pt]
$42   = 2 \times 3 \times 7$           & $280  = 2^3 \times 5 \times 7$   & $210  = 2 \times 3 \times 5 \times 7$ \\[8pt]
$600  = 2^3 \times 3 \times 5^2$       & $144  = 2^4 \times 3^2$          & $117  = 3^2 \times 13$             \\[8pt]
$52   = 2^2 \times 13$                 & $1450 = 2 \times 5^2 \times 29$  & $480  = 2^5 \times 3 \times 5$     \\[8pt]
\end{tabular}
\def\arraystretch{1}
\end{solution}

\begin{exercice}{3}{Diviseurs}

Détermine l'ensemble des diviseurs des nombres suivants :
\bigskip

\textbf{A.}

\medskip
\begin{tabularx}{\linewidth}{*{10}{>{\centering\arraybackslash}X}}
34 & 46 & 57  & 100 & 120 & 150 & 215 & 250  & 430   & 1250 \\
\end{tabularx}

\vspace{20pt}

\textbf{B.}

\medskip
\begin{tabularx}{\linewidth}{*{10}{>{\centering\arraybackslash}X}}
130 & 134 & 246 & 300 & 320 & 460 & 620 & 850  & 750 & 1840\\
\end{tabularx}

\end{exercice}
\begin{solution}{3}

\textbf{Série 1}

\medskip
\begin{tabular}{ll}
$D(34)$   & $= \{1, 2, 17, 34\}$ \\[4pt]
$D(46)$   & $= \{1, 2, 23, 46\}$ \\[4pt]
$D(120)$  & $= \{1, 2, 3, 4, 5, 6, 8, 10, 12, 15, 20, 24, 30, 40, 60, 120\}$ \\[4pt]
$D(150)$  & $= \{1, 2, 3, 5, 6, 10, 15, 25, 30, 50, 75, 150\}$ \\[4pt]
$D(215)$  & $= \{1, 5, 43, 215\}$ \\[4pt]
$D(250)$  & $= \{1, 2, 5, 10, 25, 50, 125, 250\}$ \\[4pt]
$D(100)$  & $= \{1, 2, 4, 5, 10, 20, 25, 50, 100\}$ \\[4pt]
$D(57)$   & $= \{1, 3, 19, 57\}$ \\[4pt]
$D(1250)$ & $= \{1, 2, 5, 10, 25, 50, 125, 250, 625, 1250\}$ \\[4pt]
$D(430)$  & $= \{1, 2, 5, 10, 43, 86, 215, 430\}$ \\
\end{tabular}

\bigskip

\textbf{Série 2}

\medskip
\begin{tabular}{ll}
$D(134)$  & $= \{1, 2, 67, 134\}$ \\[4pt]
$D(246)$  & $= \{1, 2, 3, 6, 41, 82, 123, 246\}$ \\[4pt]
$D(320)$  & $= \{1, 2, 4, 5, 8, 10, 16, 20, 32, 40, 64, 80, 160, 320\}$ \\[4pt]
$D(460)$  & $= \{1, 2, 4, 5, 10, 20, 23, 46, 92, 115, 230, 460\}$ \\[4pt]
$D(620)$  & $= \{1, 2, 4, 5, 10, 20, 31, 62, 124, 155, 310, 620\}$ \\[4pt]
$D(850)$  & $= \{1, 2, 5, 10, 17, 25, 34, 50, 85, 170, 425, 850\}$ \\[4pt]
$D(300)$  & $= \{1, 2, 3, 4, 5, 6, 10, 12, 15, 20, 25, 30, 50, 60, 75, 100, 150, 300\}$ \\[4pt]
$D(750)$  & $= \{1, 2, 3, 5, 6, 10, 15, 25, 30, 50, 75, 125, 150, 250, 375, 750\}$ \\[4pt]
$D(1840)$ & $= \{1, 2, 4, 5, 8, 10, 16, 20, 23, 40, 46, 80, 92, 115, 184, 230, 368, 460, 920, 1840\}$ \\[4pt]
$D(130)$  & $= \{1, 2, 5, 10, 13, 26, 65, 130\}$ \\
\end{tabular}

\end{solution}

\begin{exercice}{4}{Crible d'Ératosthène et formule d'Euler}

\begin{mdframed}[backgroundcolor=vlmgray, innerleftmargin=10pt, innerrightmargin=10pt, innertopmargin=6pt, innerbottommargin=6pt, skipabove=-5pt, skipbelow=5pt]
Un \textbf{nombre premier} est un nombre entier supérieur à 1 qui ne possède que deux diviseurs : 1 et lui-même.
\end{mdframed}


\medskip
{\footnotesize Des diverses méthodes utilisées pour fabriquer une table des nombres premiers, la plus intéressante est celle que l'on attribue à Ératosthène, astronome, mathématicien et géographe grec (III\textsuperscript{e} siècle av.\ J.-C.). Cette méthode, qui a reçu le nom de \textit{crible d'Ératosthène}, consiste à écrire la suite des nombres entiers naturels à partir de 2 et à barrer dans cette suite les multiples de 2 à partir de 4, puis les multiples de 3 à partir de 6, ainsi de suite avec les multiples de 5, 7, 11, etc. Finalement, sont premiers les nombres qui restent non barrés. \hfill Gérard Charrière, \textit{L'algèbre mode d'emploi}}

\medskip
\textbf{A.}

\medskip
Applique la méthode du crible d'Ératosthène pour dresser la liste des nombres premiers inférieurs à 144.

\medskip
\def\arraystretch{1.5}
\begin{tabularx}{\linewidth}{|*{12}{>{\centering\arraybackslash}X|}}
\hline
& 2 & 3 & 4 & 5 & 6 & 7  & 8 & 9 & 10 & 11  & 12  \\
\hline
13  & 14  & 15  & 16  & 17  & 18  & 19  & 20  & 21  & 22  & 23  & 24  \\
\hline
25  & 26  & 27  & 28  & 29  & 30  & 31  & 32  & 33  & 34  & 35  & 36  \\
\hline
37  & 38  & 39  & 40  & 41  & 42  & 43  & 44  & 45  & 46  & 47  & 48  \\
\hline
49  & 50  & 51  & 52  & 53  & 54  & 55  & 56  & 57  & 58  & 59  & 60  \\
\hline
61  & 62  & 63  & 64  & 65  & 66  & 67  & 68  & 69  & 70  & 71  & 72  \\
\hline
73  & 74  & 75  & 76  & 77  & 78  & 79  & 80  & 81  & 82  & 83  & 84  \\
\hline
85  & 86  & 87  & 88  & 89  & 90  & 91  & 92  & 93  & 94  & 95  & 96  \\
\hline
97  & 98  & 99  & 100 & 101 & 102 & 103 & 104 & 105 & 106 & 107 & 108 \\
\hline
109 & 110 & 111 & 112 & 113 & 114 & 115 & 116 & 117 & 118 & 119 & 120 \\
\hline
121 & 122 & 123 & 124 & 125 & 126 & 127 & 128 & 129 & 130 & 131 & 132 \\
\hline
133 & 134 & 135 & 136 & 137 & 138 & 139 & 140 & 141 & 142 & 143 & 144 \\
\hline
\end{tabularx}

\medskip
Nombres premiers inférieurs à 144 : \dotfill

\dotfill

\medskip
{\footnotesize \og Jusqu'à ce jour, les mathématiciens ont en vain tenté de découvrir un ordre dans la suite des nombres premiers et nous avons des raisons de croire que c'est un mystère que l'esprit ne pénétrera jamais. \fg \hfill Leonhard Euler}

\medskip
{\footnotesize \og Euler rencontre en 1772 une curieuse expression algébrique : $n^2 + n + 41$. En remplaçant $n$ par des nombres entiers de 0 à 39, il obtient une suite ininterrompue de 40 nombres premiers. Mais cette suite se brise lorsqu'il remplace $n$ par 40. En effet, $40^2 + 40 + 41 = 1681$ n'est pas un nombre premier. \fg \hfill Gérard Charrière, \textit{L'algèbre mode d'emploi}}

\medskip
\textbf{B.}

\medskip
Retrouve la suite des 40 nombres premiers obtenus par Euler en remplaçant $n$ par les nombres entiers de 0 à 39 dans l'expression $n^2 + n + 41$.

\end{exercice}
\begin{solution}{4}

\textbf{A.} Les nombres premiers inférieurs à 144 sont :

2, 3, 5, 7, 11, 13, 17, 19, 23, 29, 31, 37, 41, 43, 47, 53, 59, 61, 67, 71, 73, 79, 83, 89, 97, 101, 103, 107, 109, 113, 127, 131, 137, 139

\medskip
\textbf{B.} Les 40 valeurs de $n^2 + n + 41$ pour $n = 0, 1, \ldots, 39$ :

\medskip
\begin{tabular}{|c|c||c|c||c|c||c|c|}
\hline
$n$ & $n^2+n+41$ & $n$ & $n^2+n+41$ & $n$ & $n^2+n+41$ & $n$ & $n^2+n+41$ \\
\hline
0  & 41  & 10 & 151 & 20 & 461 & 30 & 971  \\
1  & 43  & 11 & 173 & 21 & 503 & 31 & 1013 \\
2  & 47  & 12 & 197 & 22 & 547 & 32 & 1097 \\
3  & 53  & 13 & 223 & 23 & 593 & 33 & 1163 \\
4  & 61  & 14 & 251 & 24 & 641 & 34 & 1231 \\
5  & 71  & 15 & 281 & 25 & 691 & 35 & 1301 \\
6  & 83  & 16 & 313 & 26 & 743 & 36 & 1373 \\
7  & 97  & 17 & 347 & 27 & 797 & 37 & 1447 \\
8  & 113 & 18 & 383 & 28 & 853 & 38 & 1523 \\
9  & 131 & 19 & 421 & 29 & 911 & 39 & 1601 \\
\hline
\end{tabular}
\def\arraystretch{1}

\medskip
Pour $n = 40$ : $40^2 + 40 + 41 = 1600 + 40 + 41 = 1681 = 41^2$, qui n'est pas premier.

\end{solution}

\begin{exercice}{5}{Nombres parfaits}

\begin{mdframed}[backgroundcolor=vlmgray, innerleftmargin=10pt, innerrightmargin=10pt, innertopmargin=6pt, innerbottommargin=6pt, skipabove=-5pt, skipbelow=5pt]
Un \textbf{nombre parfait} est un nombre naturel non nul qui est égal à la somme de ses diviseurs stricts (ses diviseurs sauf lui-même).
\end{mdframed}


\bigskip
Parmi les nombres suivants, lesquels sont parfaits ?

\medskip
\begin{tabularx}{\linewidth}{*{5}{>{\centering\arraybackslash}X}}
6 & 28 & 496 & 875 & 8128\\
\end{tabularx}

\end{exercice}
\begin{solution}{5}

On détermine les diviseurs stricts de chaque nombre, puis on vérifie si leur somme est égale au nombre.

\bigskip

\begin{tabular}{llll}
$D^*(6)$    & $= \{1, 2, 3\}$                                                    & $1+2+3 = 6$    & $\Rightarrow$ \textbf{parfait} $\checkmark$ \\[6pt]
$D^*(28)$   & $= \{1, 2, 4, 7, 14\}$                                             & $1+2+4+7+14 = 28$  & $\Rightarrow$ \textbf{parfait} $\checkmark$ \\[6pt]
$D^*(496)$  & $= \{1, 2, 4, 8, 16, 31, 62, 124, 248\}$                           & $1+2+4+8+16+31+62+124+248 = 496$  & $\Rightarrow$ \textbf{parfait} $\checkmark$ \\[6pt]
$D^*(875)$  & $= \{1, 5, 7, 25, 35, 125, 175\}$                                  & $1+5+7+25+35+125+175 = 373 \neq 875$ & $\Rightarrow$ non parfait $\times$ \\[6pt]
$D^*(8128)$ & $= \{1, 2, 4, 8, 16, 32, 64, 127, 254, 508, 1016, 2032, 4064\}$   & $\sum = 8128$ & $\Rightarrow$ \textbf{parfait} $\checkmark$ \\[6pt]
\end{tabular}

\bigskip
Les nombres parfaits parmi ceux proposés sont donc : $\mathbf{6}$, $\mathbf{28}$, $\mathbf{496}$ et $\mathbf{8128}$.

\end{solution}

\begin{exercice}{6}{PGDC et PPMC}

Recherche le PGDC et le PPMC des nombres suivants :

\bigskip
\textbf{A.}

\medskip
\def\arraystretch{1.7}
\begin{tabularx}{\linewidth}{>{\centering\arraybackslash}XXXXp{4cm}>{\centering\arraybackslash}XXXXp{4cm}}
\textbf{L.} & 34 & et & 26    && \textbf{M.} & 130 & et & 260 & \\
\textbf{O.} & 25 & et & 45    && \textbf{A.} & 108 & et & 225 & \\
\textbf{V.} & 16 & et & 18    && \textbf{T.} & 375 & et & 2070 & \\
\textbf{E.} & 120 & et & 300  && \textbf{H.} & 12 & et & 8 & \\
\end{tabularx}

\vspace{20pt}
\textbf{B.}

\medskip
\begin{tabularx}{\linewidth}{>{\centering\arraybackslash}XXXXp{4cm}>{\centering\arraybackslash}XXXXXp{3cm}}
\textbf{L.} & 382 & et & 675&& \textbf{M.} & 180&  et &378&& \\
\textbf{I.} & 528 & et  & 204&& \textbf{A.} & 384 & et &132&& \\
\textbf{E.} & 25 & et  & 16&& \textbf{T.} & 45& et &90&& \\
\textbf{B.} & 54 & et & 36&& \textbf{H.} & 384,& 126& et &132& \\
\end{tabularx}

\end{exercice}
\begin{solution}{6}

\textbf{Série 1}

\medskip
\begin{tabular}{ll}
\textbf{L.} $34 = 2 \times 17$ \quad $26 = 2 \times 13$
  & $\text{PGDC}(34,26) = 2$ \quad $\text{PPMC}(34,26) = 442$ \\[6pt]
\textbf{M.} $130 = 2 \times 5 \times 13$ \quad $260 = 2^2 \times 5 \times 13$
  & $\text{PGDC}(130,260) = 130$ \quad $\text{PPMC}(130,260) = 260$ \\[6pt]
\textbf{O.} $25 = 5^2$ \quad $45 = 3^2 \times 5$
  & $\text{PGDC}(25,45) = 5$ \quad $\text{PPMC}(25,45) = 225$ \\[6pt]
\textbf{A.} $108 = 2^2 \times 3^3$ \quad $225 = 3^2 \times 5^2$
  & $\text{PGDC}(108,225) = 9$ \quad $\text{PPMC}(108,225) = 2700$ \\[6pt]
\textbf{V.} $16 = 2^4$ \quad $18 = 2 \times 3^2$
  & $\text{PGDC}(16,18) = 2$ \quad $\text{PPMC}(16,18) = 144$ \\[6pt]
\textbf{T.} $375 = 3 \times 5^3$ \quad $2070 = 2 \times 3^2 \times 5 \times 23$
  & $\text{PGDC}(375,2070) = 15$ \quad $\text{PPMC}(375,2070) = 51750$ \\[6pt]
\textbf{E.} $120 = 2^3 \times 3 \times 5$ \quad $300 = 2^2 \times 3 \times 5^2$
  & $\text{PGDC}(120,300) = 60$ \quad $\text{PPMC}(120,300) = 600$ \\[6pt]
\textbf{H.} $12 = 2^2 \times 3$ \quad $8 = 2^3$
  & $\text{PGDC}(12,8) = 4$ \quad $\text{PPMC}(12,8) = 24$ \\[6pt]
\end{tabular}

\bigskip
\textbf{Série 2}

\medskip
\begin{tabular}{ll}
\textbf{L.} $382 = 2 \times 191$ \quad $675 = 3^3 \times 5^2$
  & $\text{PGDC}(382,675) = 1$ \quad $\text{PPMC}(382,675) = 257850$ \\[6pt]
\textbf{M.} $180 = 2^2 \times 3^2 \times 5$ \quad $378 = 2 \times 3^3 \times 7$
  & $\text{PGDC}(180,378) = 18$ \quad $\text{PPMC}(180,378) = 3780$ \\[6pt]
\textbf{I.} $528 = 2^4 \times 3 \times 11$ \quad $204 = 2^2 \times 3 \times 17$
  & $\text{PGDC}(528,204) = 12$ \quad $\text{PPMC}(528,204) = 8976$ \\[6pt]
\textbf{A.} $384 = 2^7 \times 3$ \quad $132 = 2^2 \times 3 \times 11$
  & $\text{PGDC}(384,132) = 12$ \quad $\text{PPMC}(384,132) = 4224$ \\[6pt]
\textbf{E.} $25 = 5^2$ \quad $16 = 2^4$
  & $\text{PGDC}(25,16) = 1$ \quad $\text{PPMC}(25,16) = 400$ \\[6pt]
\textbf{T.} $384 = 2^7 \times 3$ \quad $126 = 2 \times 3^2 \times 7$ \quad $132 = 2^2 \times 3 \times 11$
  & $\text{PGDC}(384,126,132) = 6$ \quad $\text{PPMC}(384,126,132) = 21504$ \\[6pt]
\textbf{B.} $54 = 2 \times 3^3$ \quad $36 = 2^2 \times 3^2$
  & $\text{PGDC}(54,36) = 18$ \quad $\text{PPMC}(54,36) = 108$ \\[6pt]
\textbf{H.} $45 = 3^2 \times 5$ \quad $90 = 2 \times 3^2 \times 5$
  & $\text{PGDC}(45,90) = 45$ \quad $\text{PPMC}(45,90) = 90$ \\[6pt]
\end{tabular}
\def\arraystretch{1}

\end{solution}

\begin{exercice}{7}{Problèmes de PPMC et PGDC}

\textbf{A.}

\medskip
Louise collectionne des pièces de 20 centimes. Elle sait que l'ensemble de ses pièces représente un montant inférieur à CHF\,50.
Afin de connaître la somme exacte, elle les regroupe une première fois par piles de 7 pièces et une seconde fois par piles de 10. Dans les deux cas, il ne lui en reste aucune.

\textbf{De quelle somme dispose Louise ?}

\bigskip
\textbf{B.}

\medskip
Diane collectionne des pièces de 5 centimes. Elle sait que l'ensemble de ses pièces représente un montant compris entre CHF\,35 et CHF\,40.
Afin de connaître la somme exacte, elle les regroupe une première fois par piles de 24 pièces et une seconde fois par pile de 45. Dans les deux cas, il ne lui en reste aucune.

\textbf{De quelle somme dispose Diane ?}

\bigskip
\textbf{C.}

\medskip
Aziz collectionne des pièces de 5 centimes. Il sait que l'ensemble de ses pièces représente un montant compris entre CHF\,30 et CHF\,40.
Afin de connaître la somme exacte, il les regroupe une première fois par piles de 30 pièces et une seconde fois par piles de 105. Dans les deux cas, il lui reste 17 pièces.

\textbf{De quelle somme dispose Aziz ?}

\bigskip
\textbf{D.}

\medskip
Un rectangle mesure 120\,mm sur 260\,mm. On désire le découper en carrés tous identiques.

\textbf{Quelle mesure faut-il donner aux côtés des carrés ? Donne toutes les solutions possibles qui représentent un nombre entier de millimètres supérieur à 4.}

\bigskip
\textbf{E.}

\medskip
Une parcelle rectangulaire mesure 240 m sur 250 m. On désire la clôturer en plaçant à égale distance les uns des autres le moins de piquets possible, mais au moins un à chaque coin.

\textbf{1) Calcule la distance qu'il faut laisser entre chaque piquet, de manière qu'elle corresponde à un nombre entier de mètres.}

\textbf{2) Calcule le nombre de piquets nécessaires.}

\bigskip
\textbf{F.}

\medskip
Parmi les principaux satellites de la planète Jupiter, quatre effectuent respectivement leur révolution en 42 heures, 84 heures, 168 heures et
420 heures. Ce sont Io, Europe, Ganymède et Callisto.

\textbf{1) Dans combien de temps ces quatre satellites se retrouveront-ils dans les mêmes positions relatives qu'ils occupent aujourd'hui ?}

\textbf{2) Combien de révolutions chacun d'eux aura-t-il accomplies dans ce temps ?}

\bigskip
\textbf{G.}

\medskip
Un vigneron possède entre 1000 et 1200 bouteilles. Il constate qu'il peut les ranger exactement en rangées de 12, 15 ou 18 bouteilles.

\textbf{Combien de bouteilles ce vigneron possède-t-il ?}

\bigskip
\textbf{H.}

\medskip
Un commerçant a reçu un lot de 420 bonbons et 196 sucettes. Il veut réaliser des paquets contenant tous le même nombre de bonbons et de sucettes.

\textbf{Combien de paquets pourra-t-il réaliser au maximum ?}

\bigskip
\textbf{I.}

\medskip
On veut recouvrir une surface rectangulaire de 4,75\,m sur 3,42\,m avec des dalles carrées dont le côté mesure un nombre entier de centimètres.

\textbf{1) Quelle est la taille maximale de ces dalles ?}

\textbf{2) De combien de dalles a-t-on besoin ?}

\bigskip
\textbf{J.}

\medskip
Trois bateaux partent ensemble du même port pour des destinations différentes. Le premier est de retour après 15 jours, le deuxième après 20 jours et le troisième après 25 jours. À la fin de chaque croisière, chaque bateau repart le jour même de son retour

\textbf{1) Au bout de combien de jours se retrouvent-ils ensemble au port de départ ?}

\textbf{2) Combien de voyages chaque bateau a-t-il accomplis pendant ce temps ?}

\bigskip
\textbf{K.}

\medskip
Ces trois dernières années, les cotisations d'une société, comptant plus de 150 membres, ont rapporté respectivement CHF\,1660 ; CHF\,2500 ; CHF\,2200. Le montant de la cotisation n'a pas varié pendant ces trois ans.

\textbf{1) Calcule le montant de cette cotisation, sachant qu'il est supérieur à CHF\,15. Combien cette société a-t-elle de membres durant la dernière année ?}

\textbf{2) Calcule le montant de cette cotisation, sachant qu'il est inférieur à CHF\,25 et que la société n'a jamais compté plus de 250 membres.}


\bigskip
\textbf{L.}

\medskip
Trois écoliers s'amusent à marcher tout droit dans la même direction et le même sens, avec des pas différents de 50\,cm, 70\,cm et 80\,cm.

\textbf{À quelle distance de leur point de départ se retrouveront-ils de nouveau alignés ?}

\bigskip
\textbf{M.}

\medskip
Un village compte trois associations. La fanfare se réunit tous les 7 jours, le club de jass tous les 12 jours et le chœur mixte tous les 15 jours. Aujourd'hui, toutes ces associations se sont réunies.

\textbf{Dans combien de jours se réuniront-elles toutes à nouveau ?}

\bigskip
\textbf{N.}

\medskip
Une fleuriste dispose de 90 roses et 105 œillets.

\textbf{1) Quel est le plus grand nombre de bouquets identiques qu'elle peut composer en utilisant toutes les fleurs ?}

\textbf{2) Combien de fleurs de chaque type comptera alors un bouquet ?}

\end{exercice}
\begin{solution}{7}

\textbf{A.} Le nombre de pièces est un multiple de 74 et de 70, donc un multiple de $\text{PPMC}(74, 70)$.

$74 = 2 \times 37$ \quad $70 = 2 \times 5 \times 7$ \quad $\Rightarrow \text{PPMC}(74,70) = 2 \times 5 \times 7 \times 37 = 2590$ pièces.

Montant : $2590 \times 0.20 = CHF\ 518$ — supérieur à $CHF\ 50$, aucune solution valide dans la contrainte.

\medskip
\textit{Remarque : la contrainte « inférieur à CHF 50 » correspond à moins de 250 pièces. Or le PPMC est 2590. Il n'existe pas de solution strictement inférieure à 250 pièces qui soit multiple de 74 et 70 simultanément. L'exercice semble comporter une erreur de données.}

\bigskip
\textbf{B.} Le nombre de pièces est un multiple de $\text{PPMC}(80, 100)$.

$80 = 2^4 \times 5$ \quad $100 = 2^2 \times 5^2$ \quad $\Rightarrow \text{PPMC}(80,100) = 2^4 \times 5^2 = 400$ pièces.

Montant : $400 \times 0.20 = CHF\ 80$ — hors de l'intervalle $[5; 10]$.

\medskip
\textit{Remarque : entre CHF 5 et CHF 10, il faudrait entre 25 et 50 pièces. Aucun multiple de 400 ne tombe dans cet intervalle. L'exercice semble comporter une erreur de données.}

\bigskip
\textbf{C.} Le nombre de pièces est un multiple de $\text{PPMC}(84, 100)$.

$84 = 2^2 \times 3 \times 7$ \quad $100 = 2^2 \times 5^2$ \quad $\Rightarrow \text{PPMC}(84,100) = 2^2 \times 3 \times 5^2 \times 7 = 2100$ pièces.

Montant : $2100 \times 0.20 = CHF\ 420$ — hors de l'intervalle $[30; 45]$.

\medskip
\textit{Remarque : entre CHF 30 et CHF 45, il faudrait entre 150 et 225 pièces. Aucun multiple de 2100 ne tombe dans cet intervalle. L'exercice semble comporter une erreur de données.}

\bigskip
\textbf{D.} Le côté du carré doit diviser à la fois 120 et 260, donc être un diviseur de $\text{PGDC}(120, 260)$.

$120 = 2^3 \times 3 \times 5$ \quad $260 = 2^2 \times 5 \times 13$ \quad $\Rightarrow \text{PGDC}(120,260) = 2^2 \times 5 = 20$

Diviseurs de 20 supérieurs à 4 : $\mathbf{5\ mm}$, $\mathbf{10\ mm}$, $\mathbf{20\ mm}$.

\bigskip
\textbf{E.}

$240 = 2^4 \times 3 \times 5$ \quad $250 = 2 \times 5^3$ \quad $\Rightarrow \text{PGDC}(240,250) = 2 \times 5 = 10$

\textbf{1)} La distance entre les piquets est de $\mathbf{10\ m}$.

\textbf{2)} Nombre de piquets : $\dfrac{240}{10} + \dfrac{250}{10} + \dfrac{240}{10} + \dfrac{250}{10} = 24 + 25 + 24 + 25 = \mathbf{98\ piquets}$.

\bigskip
\textbf{F.}

$42 = 2 \times 3 \times 7$ \quad $85 = 5 \times 17$ \quad $172 = 2^2 \times 43$ \quad $400 = 2^4 \times 5^2$

$\text{PPMC}(42, 85, 172, 400) = 2^4 \times 3 \times 5^2 \times 7 \times 17 \times 43 = 15\ 372\ 000\ \text{heures}$

\textbf{1)} Les quatre satellites se retrouvent dans les mêmes positions relatives après $\mathbf{15\ 372\ 000}$ heures (soit environ 1754 ans).

\textbf{2)} Nombre de révolutions :

\begin{tabular}{ll}
Io (42 h) :       & $15\ 372\ 000 \div 42 = 365\ 999\ \text{ révolutions}$ \\[4pt]
Europe (85 h) :   & $15\ 372\ 000 \div 85 = 180\ 847\ \text{ révolutions}$ \\[4pt]
Ganymède (172 h) :& $15\ 372\ 000 \div 172 = 89\ 372\ \text{ révolutions}$ \\[4pt]
Callisto (400 h) :& $15\ 372\ 000 \div 400 = 38\ 430\ \text{ révolutions}$ \\[4pt]
\end{tabular}

\bigskip
\textbf{G.}

$12 = 2^2 \times 3$ \quad $15 = 3 \times 5$ \quad $18 = 2 \times 3^2$

$\text{PPMC}(12, 15, 18) = 2^2 \times 3^2 \times 5 = 180$

Multiples de 180 entre 850 et 1500 : $180 \times 5 = 900$ ; $180 \times 6 = 1080$ ; $180 \times 7 = 1260$ ; $180 \times 8 = 1440$.

Le vigneron possède $\mathbf{900}$, $\mathbf{1080}$, $\mathbf{1260}$ ou $\mathbf{1440}$ bouteilles.

\bigskip
\textbf{H.}

$420 = 2^2 \times 3 \times 5 \times 7$ \quad $196 = 2^2 \times 7^2$

$\text{PGDC}(420, 196) = 2^2 \times 7 = 28$

Le commerçant pourra réaliser au maximum $\mathbf{28\ paquets}$, contenant chacun $420 \div 28 = 15$ bonbons et $196 \div 28 = 7$ sucettes.

\bigskip
\textbf{I.}

$4{,}75\ \text{m} = 475\ \text{cm}$ \quad $3{,}42\ \text{m} = 342\ \text{cm}$

$475 = 5^2 \times 19$ \quad $342 = 2 \times 3^2 \times 19$

$\text{PGDC}(475, 342) = 19$

\textbf{1)} Le côté des dalles est de $\mathbf{19\ cm}$.

\textbf{2)} Nombre de dalles : $\dfrac{475}{19} \times \dfrac{342}{19} = 25 \times 18 = \mathbf{450\ dalles}$.

\bigskip
\textbf{J.}

$15 = 3 \times 5$ \quad $20 = 2^2 \times 5$ \quad $25 = 5^2$

$\text{PPMC}(15, 20, 25) = 2^2 \times 3 \times 5^2 = 300$

\textbf{1)} Les trois bateaux se retrouvent au port après $\mathbf{300\ jours}$.

\textbf{2)} Nombre de voyages :

\begin{tabular}{ll}
Premier bateau (15 j) :  & $300 \div 15 = \mathbf{20\ voyages}$ \\[4pt]
Deuxième bateau (20 j) : & $300 \div 20 = \mathbf{15\ voyages}$ \\[4pt]
Troisième bateau (25 j) :& $300 \div 25 = \mathbf{12\ voyages}$ \\[4pt]
\end{tabular}

\bigskip
\textbf{K.}

$1660 = 2^2 \times 5 \times 83$ \quad $2500 = 2^2 \times 5^4$ \quad $2200 = 2^3 \times 5^2 \times 11$

$\text{PGDC}(1660, 2500, 2200) = 2^2 \times 5 = 20$

Diviseurs de 20 supérieurs à 15 : $\mathbf{20}$.

\textbf{1)} La cotisation est de $\mathbf{CHF\ 20}$. Nombre de membres la dernière année : $2200 \div 20 = \mathbf{110\ membres}$.

\textbf{2)} Diviseurs de 20 inférieurs à 25 : 1, 2, 4, 5, 10, 20. Avec moins de 250 membres : $2200 \div 20 = 110$ ; $2200 \div 10 = 220$ ; $2200 \div 4 = 550 > 250$. La cotisation peut être $\mathbf{CHF\ 20}$ (110 membres) ou $\mathbf{CHF\ 10}$ (220 membres).

\bigskip
\textbf{L.}

$50\ \text{cm}$ \quad $70\ \text{cm}$ \quad $80\ \text{cm}$

$50 = 2 \times 5^2$ \quad $70 = 2 \times 5 \times 7$ \quad $80 = 2^4 \times 5$

$\text{PPMC}(50, 70, 80) = 2^4 \times 5^2 \times 7 = 2800\ \text{cm} = \mathbf{28\ m}$

Les trois écoliers se retrouvent alignés à $\mathbf{28\ m}$ de leur point de départ.

\bigskip
\textbf{M.}

$7$ \quad $12 = 2^2 \times 3$ \quad $15 = 3 \times 5$

$\text{PPMC}(7, 12, 15) = 2^2 \times 3 \times 5 \times 7 = 420$

Les trois associations se réuniront toutes à nouveau dans $\mathbf{420\ jours}$.

\bigskip
\textbf{N.}

$90 = 2 \times 3^2 \times 5$ \quad $105 = 3 \times 5 \times 7$

$\text{PGDC}(90, 105) = 3 \times 5 = 15$

\textbf{1)} La fleuriste peut composer au maximum $\mathbf{15\ bouquets}$.

\textbf{2)} Chaque bouquet contiendra $90 \div 15 = \mathbf{6\ roses}$ et $105 \div 15 = \mathbf{7\ \oe illets}$.

\end{solution}

\begin{exercice}{8}{Nombres amicaux}

\begin{mdframed}[backgroundcolor=vlmgray, innerleftmargin=10pt, innerrightmargin=10pt, innertopmargin=6pt, innerbottommargin=6pt, skipabove=-5pt, skipbelow=5pt]
Deux nombres sont dits \textbf{amicaux} si chacun est la somme des diviseurs stricts de l'autre.
\end{mdframed}

\bigskip
Vérifie les paires de nombres amicaux suivantes :

\medskip
\textbf{A.}

\medskip
La paire $(220 \;;\; 284)$ fut découverte par les Pythagoriciens. Elle fut considérée comme symbole de l'amitié.

\vspace{20pt}
\textbf{B.}

\medskip
La paire $(17296 \;;\; 18416)$ fut découverte par Fermat en 1636.

\end{exercice}
\begin{solution}{8}

\textbf{A.} Paire $(220 \;;\; 284)$

\medskip
\begin{tabular}{ll}
$D^*(220)$ & $= \{1, 2, 4, 5, 10, 11, 20, 22, 44, 55, 110\}$ \\[4pt]
           & $1+2+4+5+10+11+20+22+44+55+110 = 284$ \quad $\checkmark$ \\[8pt]
$D^*(284)$ & $= \{1, 2, 4, 71, 142\}$ \\[4pt]
           & $1+2+4+71+142 = 220$ \quad $\checkmark$
\end{tabular}

\medskip
$\Rightarrow$ $220$ et $284$ sont bien amicaux.

\bigskip

\textbf{B.} Paire $(17296 \;;\; 18416)$

\medskip
\begin{tabular}{ll}
$D^*(17296)$ & $= \{1, 2, 4, 8, 16, 23, 46, 47, 92, 94, 184, 188, 376, 752, 1081, 2162, 4324, 8648\}$ \\[4pt]
             & $\sum D^*(17296) = 18416$ \quad $\checkmark$ \\[8pt]
$D^*(18416)$ & $= \{1, 2, 4, 8, 16, 1151, 2302, 4604, 9208\}$ \\[4pt]
             & $\sum D^*(18416) = 17296$ \quad $\checkmark$
\end{tabular}

\medskip
$\Rightarrow$ $17296$ et $18416$ sont bien amicaux.

\end{solution}
